% page 589 du fac-simile (image p-588 du scan)
% bloc 180.3 x 246.3 mm ; marge gauche 21.7 mm
\pgc{0.000mm}{0.000mm}{
\l{}
\l{\cel{6}\sou{tilio}. (\sou{bot.}) Arboro di qua uzesas : la floro, por facar infu-}
\l{\cel{9}zuri kalmigiva, e la kortico, por facar puteo-kordegi.- L}
\l{\cel{9}\sou{tilia}. - FISL.}
\l{}
\l{\cel{6}\sou{timar}. (\sou{trans}.) I. Tendencar evitor (ulu, ulo) (de qua onu}
\l{\cel{9}expektas malajo). - II. Judikar ulo (regretinda, desoportuna)}
\l{\cel{9}kom eventonta plu o min balde. - EFIS.}
\l{}
\l{\cel{6}\sou{timbalo}. Mi-globo de kupro o latuno, kovrita per felo tensita}
\l{\cel{9}quan onu klemas o desklemas per skrubi, por akordigar lu,}
\l{\cel{9}e sur qua onu frapas per vergeti ek ligno harda o tegita ye}
\l{\cel{9}felo-peco segun la sonoreso-grado quan onu deziras obtenor -}
\l{\cel{9}muzik-instrumento qua uzesas en kavalrio ed en orkestro.-EFIS.}
\l{}
\l{\cel{6}timbro (+tembro). Kaloto de metalo qua, frapate per marteleto,}
\l{\cel{9}produktas sonuno pasable duroza. - FS.}
\l{}
\l{\cel{6}+timbro (postmarko).Peceto de papero imprimuroza (ordinare kun}
\l{\cel{9}konturo dentetoza pro perforesir) quan onu glutinas sur la}
\l{\cel{9}letro-kuverti por montrar ke la afranko-preco esas pagita.}
\l{\cel{9}- Analoge pri la \textquotedbl{}timbri\textquotedbl{} fiskala, asekurala, e c.).-}
\l{\cel{9}FS.}
\l{}
\l{\cel{6}timiano. (bot.)\cel{3}Planto odoroza, ek la familio \textquotedbl{}labiei\textquotedbl{}, quan}
\l{\cel{9}onu uzas en la faco di sauci. - L.\cel{1}\sou{thymus}. - DEFIS.}
\l{}
\l{\cel{6}timida. Qua indijas audaco pro desfidar su. - EFIS.}
\l{}
\l{\cel{6}timmono.Longa stango de ligno, muntita ye la parto avana di}
\l{\cel{9}veturo, di plugilo, ed an singla latero di qua tir-bestio (ka-}
\l{\cel{9}valo, bovo) esas jungita por apogar su ad olu por tirar,}
\l{\cel{9}retenar, chanjar sua direciono, des-avancar. - FIS.}
\l{}
\l{\cel{6}timpano. I. (anat.)\cel{2}Kavajo ye qua konsistas la meza parto di}
\l{\cel{9}orelo, klozita, en la parto extera ed en la parto interna di}
\l{\cel{9}la orelo,per membrani tensita. - II. (\sou{arkitekt.})\cel{2}Spaco quan}
\l{\cel{9}cirkondas la tri kornici di frontono, qua destinesas recevor}
\l{\cel{9}basa reliefi, ornamenti. - III. (\sou{muziko}). Instrumento olima}
\l{\cel{9}di muziko, garnisita ye kordi latuna, quan onu pleas per du}
\l{\cel{9}vergeti ek ligno.\cel{2}EFIRS.}
\l{}
\l{\cel{6}timuso.\cel{2}(anat.)\cel{2}Korpo glandatra, biloba, situita dop la ster}
\l{\cel{1})\cel{7}numo, an la parto infra di la kolo. - DEFIS.}
\l{}
\l{\cel{1})\cel{4}\sou{tino}. Earelo di qua la fundo esas plu\cel{1}\sou{lar}ja kam la parto sunra}
\l{\cel{9}uzata por la transporto di aquo, vito-beri lor la rekolto,}
\l{\cel{9}e c. - DFIS.}
\l{}
\l{\cel{6}\sou{tindro}. Substanco sponjatra, extraktita de la agariko di la}
\l{\cel{9}querko (fungo kun kortico lignatra, L. ochropus fomentarius)}
\l{\cel{9}qua acendesas kontakte di cintilo. - DE.}
\l{}
\l{\cel{6}\sou{tineo.}\cel{1}(\sou{zool.}) I. Mikra insekto lepidoptera, di\cel{1}\sou{qua}\cel{1}la speco}
\l{\cel{9}maxim difuzita esas ta qua devoras la grani. - II. Mikra}
\l{\cel{9}insekto lepidoptera, noktala, di qua la larvo rodas la stofi}
\l{\cel{9}la libri, e c. - L.\cel{1}\sou{tinea}. - FISL.}
}
